%% =====================================================================
%%  B3-T-04-02 -- what B3 contributes to "The Favour Debt"
%%  -------------------------------------------------------------------
%%  LAYER A: APPEND-ONLY. Written once, then left alone. Material found
%%  only here sits in the `unique` slot and is protected by rule R10.
%% =====================================================================
%% @kind: source
%% @id: B3-T-04-02
%% @book: B3
%% @topic: T-04-02
%% @locator: Law 13, pp. 95--101; Law 12; Law 40
%% @status: distilled
%% @unique: yes
%% @integrated: yes
%% @updated: 2026-09-19

\dossier{B3}{T-04-02}{\bookshortB3}{Law 13, pp. 95--101; Law 12; Law 40}

\begin{distilled}
  \item When asking for help, appeal to the other party's self-interest and never to their mercy or their gratitude; reminding an ally of your past assistance gives them a way to ignore you.
  \item Find something in the request that benefits them and emphasise it out of proportion.
  \item Selective honesty and generosity are described as a way to lower a suspicious person's guard --- a timely gift opens a hole in the armour.
  \item The free lunch is to be despised: what is offered for nothing usually carries a trick or a hidden obligation, and paying your own way keeps you clear of gratitude, guilt and deceit.
\end{distilled}

\begin{unique}
  \item The instruction to appeal to self-interest rather than to gratitude, with the reason that a reminded debtor finds a way to ignore you.
  \item Paying full price as a deliberate strategy for staying outside other people's accounts.
\end{unique}
